\begin{enumerate}
\item Từ giả thiết, viết kết quả đo hoặc tính toán dưới dạng $m \le \overline{a} \le M$ hoặc tìm cận trên của $|\overline{a} - a|$.
\item Nếu biết $a$ là số gần đúng, độ chính xác là số $d > 0$ thỏa mãn $|\overline{a} - a| \le d$.
\item Nếu chỉ biết $m \le \overline{a} \le M$, chọn $a = \dfrac{m+M}{2}$ và độ chính xác $d = \dfrac{M-m}{2}$ để viết $\overline{a} = a \pm d$.
\end{enumerate}
Lưu ý: Số $d$ càng nhỏ thì độ chính xác càng cao; mọi số $d' > d$ cũng là một độ chính xác của $a$.
